\chapter{Dynamic Programming}
\chapterepigraph{Dynamic programming is brute force made efficient by
never solving the same subproblem twice. Every problem in the CSES DP
section is the same three questions: what is the smallest piece of state
that describes a subproblem, how does a subproblem's answer combine its
smaller subproblems, and in what order must they be filled so each
dependency is ready. Answer those three and the code is mechanical.}

\section{The Three Questions Behind Every DP}
\label{sec:cses-dp-intro}

A dynamic program is a recurrence plus memory. The art is entirely in
choosing the \emph{state}; once the state and its transition are right,
the implementation is a loop.

\begin{ideabox}[State, Transition, Order]
For every problem ask:
\begin{enumerate}
  \item \textbf{State.} What parameters uniquely identify a subproblem?
        (a prefix length, a remaining sum, a cell $(i,j)$, a bitmask of
        used items.) The number of states times the cost per state is
        the complexity.
  \item \textbf{Transition.} How is $\textit{dp}[\text{state}]$ built
        from smaller states? This is the recurrence -- usually ``try
        every last choice.''
  \item \textbf{Base case \& order.} Which states are known outright,
        and in what order must states be computed so every dependency is
        already done? (Iterative bottom-up, or memoised top-down.)
\end{enumerate}
\end{ideabox}

\subsection*{How to Recognise the DP Type}

\begin{patternbox}[CSES DP Families]
\begin{itemize}
  \item \textbf{``Count the number of ways'' / ``minimum coins''} --
        one-dimensional DP over a target value: Dice Combinations,
        Minimizing Coins, Coin Combinations I/II, Removing Digits,
        Money Sums. The key subtlety is \emph{ordered vs unordered}
        counting (permutations vs combinations), which is decided by
        loop order.
  \item \textbf{Grid / two-sequence DP} -- state is a cell $(i,j)$:
        Grid Paths, Edit Distance, Counting Tilings.
  \item \textbf{Knapsack-style} -- items with weights/values, state is
        (prefix, capacity): Book Shop, Money Sums, Two Sets II.
  \item \textbf{Interval / game DP} -- state is a subarray $[l,r]$:
        Removal Game, Rectangle Cutting.
  \item \textbf{Bitmask DP} -- state includes a subset: Elevator Rides,
        Counting Tilings (profile), Two Sets.
  \item \textbf{Digit DP} -- count numbers with a property up to $N$
        digit by digit: Counting Numbers.
  \item \textbf{Classic sequence DP} -- Increasing Subsequence (patience
        sorting), Array Description, Projects (weighted interval
        scheduling), Counting Towers.
\end{itemize}
\end{patternbox}

Each section states the problem, identifies the state and transition,
gives the recurrence and its order, a hand-traced dry run, complete
compilable C++ (using \texttt{long long} and the CSES modulus
$10^9+7$ where the problem counts things), and a complexity analysis.

\newpage

\section{Dice Combinations}
\label{sec:cses-dice-combinations}

\problemstatement{Count the number of ways to construct the sum $n$ by
throwing a die one or more times, where each throw yields $1$ to $6$ and
different \emph{orders} count as different ways. Output the count modulo
$10^9+7$.}

\begin{patternbox}
\emph{``Number of ways'' + ``order matters''} is a one-dimensional
counting DP over the target sum. Let $\textit{dp}[x]$ be the number of
ordered sequences summing to $x$; the last throw was some value
$1\ldots6$, so $\textit{dp}[x]$ sums the ways to have reached
$x-1,\ldots,x-6$ before that final throw.
\end{patternbox}

\subsection*{State and transition}

Define $\textit{dp}[x]$ as the number of ways to obtain sum $x$. The
final die throw contributes a value $k\in\{1,\ldots,6\}$, and everything
before it must sum to $x-k$. Since every different final value (and every
different earlier sequence) is a distinct ordered way, we simply add:
\[
  \textit{dp}[x]=\sum_{k=1}^{6}\textit{dp}[x-k],\qquad \textit{dp}[0]=1.
\]
The base $\textit{dp}[0]=1$ counts the single empty sequence.

\begin{ideabox}[Ordered Counting: Sum Over the Last Move]
Because order matters, we classify sequences by their \emph{last} throw.
Summing over all six possible last throws counts every ordered sequence
exactly once, with no double counting -- each sequence has a unique final
element.
\end{ideabox}

\subsection*{Algorithm}

\begin{enumerate}
  \item $\textit{dp}[0]=1$.
  \item For $x=1\ldots n$: $\textit{dp}[x]=\sum_{k=1}^{6}
        \textit{dp}[x-k]$ (skipping $x-k<0$), taken modulo $10^9+7$.
  \item Answer is $\textit{dp}[n]$.
\end{enumerate}

\subsection*{Dry run}

\dryrun{$n=3$.}

\begin{itemize}
  \item $\textit{dp}[0]=1$.
  \item $\textit{dp}[1]=\textit{dp}[0]=1$ (just ``1'').
  \item $\textit{dp}[2]=\textit{dp}[1]+\textit{dp}[0]=2$ (``1+1'',
        ``2'').
  \item $\textit{dp}[3]=\textit{dp}[2]+\textit{dp}[1]+\textit{dp}[0]
        =2+1+1=4$ (``1+1+1'',``1+2'',``2+1'',``3'').
\end{itemize}

\subsection*{C++ implementation}

\begin{lstlisting}
#include <bits/stdc++.h>
using namespace std;
const int MOD = 1e9 + 7;

int main() {
    int n;
    cin >> n;
    vector<long long> dp(n + 1, 0);
    dp[0] = 1;
    for (int x = 1; x <= n; x++)
        for (int k = 1; k <= 6; k++)
            if (x - k >= 0) dp[x] = (dp[x] + dp[x - k]) % MOD;
    cout << dp[n] << "\n";
    return 0;
}
\end{lstlisting}

\begin{complexitybox}
\textbf{Time Complexity.} $O(6n)=O(n)$ -- each of $n$ states sums six
predecessors. \textbf{Space Complexity.} $O(n)$ for the DP array (reducible
to $O(1)$ by keeping only the last six values).
\end{complexitybox}

\begin{takeawaybox}
The archetypal ordered-counting DP: \emph{sum over the last move}. Every
``count sequences where order matters'' problem in this chapter is a
variation on this single recurrence -- the only thing that changes is the
set of allowed last moves.
\end{takeawaybox}

\section{Minimizing Coins}
\label{sec:cses-minimizing-coins}

\problemstatement{Given $n$ coin values and a target sum $x$, find the
\emph{minimum} number of coins needed to make exactly $x$ (each coin may
be used any number of times), or report $-1$ if it is impossible.}

\begin{patternbox}
\emph{``Minimum coins to make a sum,'' unlimited supply}, is the
\textbf{unbounded knapsack} minimisation. Let $\textit{dp}[s]$ be the
fewest coins summing to $s$; the last coin used is some value $c$, so
$\textit{dp}[s]=1+\min_c \textit{dp}[s-c]$.
\end{patternbox}

\subsection*{State and transition}

$\textit{dp}[s]$ = minimum coins to make sum $s$, with
$\textit{dp}[0]=0$ and all others initialised to ``infinity''
(unreachable). For each $s$, try every coin $c\le s$ as the last coin:
if $s-c$ is reachable, $s$ can be reached with one more coin. Take the
best:
\[
  \textit{dp}[s]=\min_{c:\,c\le s}\big(\textit{dp}[s-c]+1\big).
\]

\begin{ideabox}[Minimisation Instead of Counting]
This is Dice Combinations with $\min+1$ replacing summation: instead of
counting all ways to make the last move, we keep only the cheapest.
Coins may repeat, so every coin is available at every sum -- that is what
makes it \emph{unbounded}.
\end{ideabox}

\subsection*{Algorithm}

\begin{enumerate}
  \item $\textit{dp}[0]=0$; $\textit{dp}[s]=\infty$ for $s\ge1$.
  \item For $s=1\ldots x$ and each coin $c\le s$: if $\textit{dp}[s-c]$
        is finite, $\textit{dp}[s]=\min(\textit{dp}[s],
        \textit{dp}[s-c]+1)$.
  \item Answer is $\textit{dp}[x]$, or $-1$ if it is still $\infty$.
\end{enumerate}

\subsection*{Dry run}

\dryrun{Coins $\{1,5,7\}$, $x=11$.}

\begin{itemize}
  \item $\textit{dp}[5]=1$ (one 5), $\textit{dp}[7]=1$,
        $\textit{dp}[10]=2$ ($5+5$).
  \item $\textit{dp}[11]$: via coin $1$ from $\textit{dp}[10]=2\to3$;
        via coin $5$ from $\textit{dp}[6]$ ($=2$, as $5{+}1$)$\to3$; via
        coin $7$ from $\textit{dp}[4]$ ($=4$)$\to5$. Minimum $=3$.
  \item Answer $3$ (e.g.\ $5+5+1$).
\end{itemize}

\subsection*{C++ implementation}

\begin{lstlisting}
#include <bits/stdc++.h>
using namespace std;

int main() {
    int n, x;
    cin >> n >> x;
    vector<int> coins(n);
    for (int& c : coins) cin >> c;

    const int INF = 1e9;
    vector<int> dp(x + 1, INF);
    dp[0] = 0;
    for (int s = 1; s <= x; s++)
        for (int c : coins)
            if (c <= s && dp[s - c] != INF)
                dp[s] = min(dp[s], dp[s - c] + 1);

    cout << (dp[x] == INF ? -1 : dp[x]) << "\n";
    return 0;
}
\end{lstlisting}

\begin{complexitybox}
\textbf{Time Complexity.} $O(nx)$ -- each of $x$ sums tries all $n$
coins. \textbf{Space Complexity.} $O(x)$ for the DP array.
\end{complexitybox}

\begin{takeawaybox}
Swap ``sum of ways'' for ``$\min$ of ways $+1$'' and a counting DP becomes
an optimisation DP. The unbounded (coins reusable) structure is why every
coin is tried at every sum, distinguishing this from the $0/1$ knapsack of
Book Shop (Section~\ref{sec:cses-book-shop}).
\end{takeawaybox}

\section{Coin Combinations I}
\label{sec:cses-coin-combinations-i}

\problemstatement{Given $n$ coins and a target $x$, count the number of
\emph{ordered} ways to make $x$ (a coin may be used many times, and
sequences differing only in order are counted separately). Output modulo
$10^9+7$.}

\begin{patternbox}
\emph{``Count ways, order matters''} with reusable coins is Dice
Combinations generalised to an arbitrary coin set. Let $\textit{dp}[s]$
count ordered sequences summing to $s$; the last coin is any $c$, so
$\textit{dp}[s]=\sum_c \textit{dp}[s-c]$.
\end{patternbox}

\subsection*{State, transition, and loop order}

$\textit{dp}[s]$ counts ordered sequences summing to $s$,
$\textit{dp}[0]=1$. The last coin used has some value $c$; the prefix
before it sums to $s-c$:
\[
  \textit{dp}[s]=\sum_{c:\,c\le s}\textit{dp}[s-c].
\]
The crucial detail is the \textbf{loop order}: the \emph{sum} $s$ runs on
the outside and coins on the inside. This way every coin is considered as
a possible \emph{last} move for each sum, so $(2\text{ then }3)$ and
$(3\text{ then }2)$ are both counted -- exactly what ``order matters''
requires.

\begin{ideabox}[Sum-Outer Loop Counts Permutations]
Putting the sum on the outer loop lets each sum choose \emph{any} coin as
its last step, generating all orderings. Reversing the loops (coins
outer) would count each multiset once -- the distinction between this
problem and Coin Combinations II (Section~\ref{sec:cses-coin-combinations-ii}).
\end{ideabox}

\subsection*{Algorithm}

\begin{enumerate}
  \item $\textit{dp}[0]=1$.
  \item For $s=1\ldots x$ (outer): for each coin $c\le s$ (inner),
        $\textit{dp}[s]\mathrel{+}=\textit{dp}[s-c]$, modulo $10^9+7$.
  \item Answer $\textit{dp}[x]$.
\end{enumerate}

\subsection*{Dry run}

\dryrun{Coins $\{2,3,5\}$, $x=9$; partial values.}

\begin{itemize}
  \item $\textit{dp}[0]=1$, $\textit{dp}[2]=1$, $\textit{dp}[3]=1$,
        $\textit{dp}[5]=\textit{dp}[3]+\textit{dp}[2]+\textit{dp}[0]=3$
        ($2{+}3$, $3{+}2$, $5$).
  \item Building up, $\textit{dp}[9]=8$ -- all ordered sequences of
        $\{2,3,5\}$ summing to $9$.
\end{itemize}

\subsection*{C++ implementation}

\begin{lstlisting}
#include <bits/stdc++.h>
using namespace std;
const int MOD = 1e9 + 7;

int main() {
    int n, x;
    cin >> n >> x;
    vector<int> coins(n);
    for (int& c : coins) cin >> c;

    vector<long long> dp(x + 1, 0);
    dp[0] = 1;
    for (int s = 1; s <= x; s++)          // sum outer -> permutations
        for (int c : coins)
            if (c <= s) dp[s] = (dp[s] + dp[s - c]) % MOD;

    cout << dp[x] << "\n";
    return 0;
}
\end{lstlisting}

\begin{complexitybox}
\textbf{Time Complexity.} $O(nx)$. \textbf{Space Complexity.} $O(x)$.
\end{complexitybox}

\begin{takeawaybox}
Ordered coin counting = sum-outer, coin-inner loops. Memorise this loop
order as the switch that toggles between counting permutations (here) and
combinations (next section). Getting it backward is the single most common
DP counting bug.
\end{takeawaybox}

\section{Coin Combinations II}
\label{sec:cses-coin-combinations-ii}

\problemstatement{Given $n$ coins and a target $x$, count the number of
\emph{distinct} ways to make $x$ where the \emph{order does not matter}
(each multiset of coins is counted once). Output modulo $10^9+7$.}

\begin{patternbox}
\emph{``Count ways, order does not matter''} is the combination count.
The trick that removes duplicate orderings is to fix a global order on
the coins and only ever extend a combination with coins ``at or after''
the current one -- realised by putting the \textbf{coin loop on the
outside}.
\end{patternbox}

\subsection*{Why coin-outer counts each multiset once}

Process coins one at a time. After incorporating the first $i$ coins,
$\textit{dp}[s]$ counts the combinations of \emph{those coins only} that
sum to $s$. Adding coin $c_i$ updates
$\textit{dp}[s]\mathrel{+}=\textit{dp}[s-c_i]$ for increasing $s$, which
appends copies of $c_i$ to combinations already built from coins
$\le i$. Because a combination is built in a fixed coin order, each
multiset is generated exactly once -- never as a reordering.

\begin{ideabox}[Coin-Outer Loop Counts Combinations]
Freezing the coin order and only extending forward eliminates permutations
of the same multiset. This is the exact mirror of Coin Combinations~I:
same array update, loops swapped. Coin-outer $\Rightarrow$ combinations;
sum-outer $\Rightarrow$ permutations.
\end{ideabox}

\subsection*{Algorithm}

\begin{enumerate}
  \item $\textit{dp}[0]=1$.
  \item For each coin $c$ (outer): for $s=c\ldots x$ (inner),
        $\textit{dp}[s]\mathrel{+}=\textit{dp}[s-c]$ modulo $10^9+7$.
  \item Answer $\textit{dp}[x]$.
\end{enumerate}

\subsection*{Dry run}

\dryrun{Coins $\{2,3,5\}$, $x=9$.}

\begin{itemize}
  \item After coin $2$: $\textit{dp}[\text{even}]=1$ for $0,2,4,6,8$.
  \item After coin $3$: combinations of $\{2,3\}$; e.g.\
        $\textit{dp}[5]=1$ ($2{+}3$), $\textit{dp}[9]$ grows.
  \item After coin $5$: final $\textit{dp}[9]=3$ -- the multisets
        $\{2,2,5\}$, $\{2,2,2,3\}$, $\{3,3,3\}$.
\end{itemize}

\subsection*{C++ implementation}

\begin{lstlisting}
#include <bits/stdc++.h>
using namespace std;
const int MOD = 1e9 + 7;

int main() {
    int n, x;
    cin >> n >> x;
    vector<int> coins(n);
    for (int& c : coins) cin >> c;

    vector<long long> dp(x + 1, 0);
    dp[0] = 1;
    for (int c : coins)                   // coin outer -> combinations
        for (int s = c; s <= x; s++)
            dp[s] = (dp[s] + dp[s - c]) % MOD;

    cout << dp[x] << "\n";
    return 0;
}
\end{lstlisting}

\begin{complexitybox}
\textbf{Time Complexity.} $O(nx)$. \textbf{Space Complexity.} $O(x)$.
\end{complexitybox}

\begin{takeawaybox}
Combinations vs permutations in coin DP is decided by \emph{which loop is
outer}, nothing else. Keep the pair Coin Combinations I and II side by
side in memory as the canonical illustration -- the same one-line update,
opposite loop nesting, opposite meaning.
\end{takeawaybox}

\section{Removing Digits}
\label{sec:cses-removing-digits}

\problemstatement{Start from an integer $n$. On each step you may
subtract any one of the digits currently appearing in the number. Find
the minimum number of steps to reduce $n$ to $0$.}

\begin{patternbox}
\emph{``Minimum steps to reduce a number to zero''} is a shortest-path
DP over values $0\ldots n$. Let $\textit{dp}[v]$ be the fewest steps from
$v$ to $0$; from $v$ you may subtract any digit of $v$, so
$\textit{dp}[v]=1+\min_{d\in\text{digits}(v)}\textit{dp}[v-d]$.
\end{patternbox}

\subsection*{State and transition}

$\textit{dp}[v]$ = minimum steps to bring $v$ down to $0$, with
$\textit{dp}[0]=0$. The only moves from $v$ subtract one of $v$'s own
decimal digits. Since a larger digit removes more in one step but the
available digits depend on $v$ itself, we must try every digit of $v$ and
take the minimum successor:
\[
  \textit{dp}[v]=1+\min_{d\,\in\,\text{digits}(v),\ d>0}\textit{dp}[v-d].
\]

\begin{ideabox}[Greedy ``Subtract the Largest Digit'' Also Works]
Because subtracting a larger digit never overshoots ($v-d\ge0$ always,
as $d\le9\le v$ for $v\ge10$, and single digits vanish in one step), a
greedy that always removes the largest digit is optimal here. The DP is
shown because it generalises and needs no proof of the greedy exchange
argument.
\end{ideabox}

\subsection*{Algorithm}

\begin{enumerate}
  \item $\textit{dp}[0]=0$.
  \item For $v=1\ldots n$: for each non-zero digit $d$ of $v$,
        $\textit{dp}[v]=\min(\textit{dp}[v], \textit{dp}[v-d]+1)$.
  \item Answer $\textit{dp}[n]$.
\end{enumerate}

\subsection*{Dry run}

\dryrun{$n=27$.}

\begin{itemize}
  \item From $27$ subtract $7\to20$; from $20$ subtract $2\to18$; then
        $8\to10$; then $1\to9$; then $9\to0$.
  \item The DP finds $\textit{dp}[27]=5$ via the value chain
        $27\to20\to18\to10\to9\to0$ (each step removes one digit).
\end{itemize}

\subsection*{C++ implementation}

\begin{lstlisting}
#include <bits/stdc++.h>
using namespace std;

int main() {
    int n;
    cin >> n;
    vector<int> dp(n + 1, 1e9);
    dp[0] = 0;
    for (int v = 1; v <= n; v++) {
        int t = v;
        while (t > 0) {
            int d = t % 10;
            if (d > 0) dp[v] = min(dp[v], dp[v - d] + 1);
            t /= 10;
        }
    }
    cout << dp[n] << "\n";
    return 0;
}
\end{lstlisting}

\begin{complexitybox}
\textbf{Time Complexity.} $O(n\log_{10} n)$ -- each value inspects its
$O(\log n)$ digits. \textbf{Space Complexity.} $O(n)$ for the DP array.
\end{complexitybox}

\begin{takeawaybox}
When moves from a state depend on the state's own structure (here, its
digits), enumerate those moves inside the transition. This ``value as a
node, allowed subtractions as edges'' framing turns a reduce-to-zero
puzzle into a one-dimensional shortest-path DP.
\end{takeawaybox}

\section{Grid Paths}
\label{sec:cses-grid-paths}

\problemstatement{Given an $n\times n$ grid where some cells are traps
(\texttt{*}) and the rest are free (\texttt{.}), count the number of
paths from the top-left to the bottom-right cell moving only right or
down, never through a trap. Output modulo $10^9+7$.}

\begin{patternbox}
\emph{``Count paths in a grid, moving right/down only''} is the textbook
2D grid DP. A cell $(i,j)$ is reached only from directly above or
directly left, so $\textit{dp}[i][j]=\textit{dp}[i-1][j]+
\textit{dp}[i][j-1]$, with traps forced to $0$.
\end{patternbox}

\subsection*{State and transition}

$\textit{dp}[i][j]$ = number of valid paths from $(0,0)$ to $(i,j)$.
Since the only moves are right and down, the last step into $(i,j)$ came
from $(i-1,j)$ (a down move) or $(i,j-1)$ (a right move). Summing those
two counts every path once. A trap cell contributes zero paths, so we
set its DP value to $0$ regardless.
\[
  \textit{dp}[i][j]=\begin{cases}0 & (i,j)\text{ is a trap}\\
  \textit{dp}[i-1][j]+\textit{dp}[i][j-1] & \text{otherwise.}\end{cases}
\]

\begin{ideabox}[Two Incoming Directions, One Sum]
Grid DP works because the move set is acyclic and monotone: you can only
increase $i$ or $j$, so processing rows top-to-bottom, columns
left-to-right, guarantees both predecessors are already computed. Traps
are handled by zeroing, which naturally propagates ``no path through
here.''
\end{ideabox}

\subsection*{Algorithm}

\begin{enumerate}
  \item $\textit{dp}[0][0]=1$ if the start is free, else $0$.
  \item For each cell in row-major order: if it is a trap set $0$; else
        add the values from above and from the left (treating
        out-of-bounds as $0$).
  \item Answer $\textit{dp}[n-1][n-1]$.
\end{enumerate}

\subsection*{Dry run}

\dryrun{$3\times3$ grid with a trap at the centre $(1,1)$.}

\begin{itemize}
  \item Row $0$: $1,1,1$. Column $0$: $1,1,1$.
  \item $(1,1)$ is a trap $\to0$. $(1,2)=\textit{dp}[0][2]+
        \textit{dp}[1][1]=1+0=1$. $(2,1)=0+1=1$.
  \item $(2,2)=\textit{dp}[1][2]+\textit{dp}[2][1]=1+1=2$.
\end{itemize}

\subsection*{C++ implementation}

\begin{lstlisting}
#include <bits/stdc++.h>
using namespace std;
const int MOD = 1e9 + 7;

int main() {
    int n;
    cin >> n;
    vector<string> g(n);
    for (auto& row : g) cin >> row;

    vector<vector<long long>> dp(n, vector<long long>(n, 0));
    for (int i = 0; i < n; i++)
        for (int j = 0; j < n; j++) {
            if (g[i][j] == '*') { dp[i][j] = 0; continue; }
            if (i == 0 && j == 0) { dp[i][j] = 1; continue; }
            long long up   = (i > 0) ? dp[i - 1][j] : 0;
            long long left = (j > 0) ? dp[i][j - 1] : 0;
            dp[i][j] = (up + left) % MOD;
        }
    cout << dp[n - 1][n - 1] << "\n";
    return 0;
}
\end{lstlisting}

\begin{complexitybox}
\textbf{Time Complexity.} $O(n^2)$ -- one pass over the grid.
\textbf{Space Complexity.} $O(n^2)$, reducible to $O(n)$ by keeping only
the previous row.
\end{complexitybox}

\begin{takeawaybox}
Grid path counting is the 2D analogue of the 1D coin sums: sum over the
directions the last move could have come from. Obstacles become zeros
that propagate automatically. This $(i,j)$-state template underlies Edit
Distance and Counting Tilings later in the chapter.
\end{takeawaybox}

\section{Book Shop}
\label{sec:cses-book-shop}

\problemstatement{You have $x$ money and there are $n$ books, each with a
price and a page count. You may buy each book at most once. Maximise the
total number of pages you can buy within your budget.}

\begin{patternbox}
\emph{``Each item at most once, maximise value under a capacity''} is the
classic \textbf{0/1 knapsack}. Price is weight, pages are value, $x$ is
capacity. Let $\textit{dp}[b]$ be the maximum pages achievable with
budget exactly $\le b$.
\end{patternbox}

\subsection*{State, transition, and the reverse inner loop}

Process books one at a time. $\textit{dp}[b]$ = best pages using books
considered so far with budget $b$. For a book with price $p$ and pages
$s$, either skip it ($\textit{dp}[b]$ unchanged) or buy it
($\textit{dp}[b-p]+s$). The decisive implementation detail: when using a
1D array, iterate $b$ \textbf{downward} from $x$ to $p$, so that
$\textit{dp}[b-p]$ still refers to a state \emph{without} the current
book -- guaranteeing each book is used at most once.
\[
  \textit{dp}[b]=\max\big(\textit{dp}[b],\ \textit{dp}[b-p]+s\big).
\]

\begin{ideabox}[Downward Loop Enforces ``At Most Once'']
In 0/1 knapsack the inner budget loop must go high-to-low. A low-to-high
loop would let $\textit{dp}[b-p]$ already include the current book,
effectively allowing it to be bought repeatedly -- which is the unbounded
knapsack (Coin Combinations II). The loop direction is the whole
difference.
\end{ideabox}

\subsection*{Algorithm}

\begin{enumerate}
  \item $\textit{dp}[b]=0$ for all $b$.
  \item For each book $(p,s)$: for $b=x$ down to $p$,
        $\textit{dp}[b]=\max(\textit{dp}[b],\textit{dp}[b-p]+s)$.
  \item Answer $\textit{dp}[x]$.
\end{enumerate}

\subsection*{Dry run}

\dryrun{Budget $x=10$; books $(\text{price},\text{pages})=(4,5),(8,12),
(5,8),(3,1)$.}

\begin{itemize}
  \item After book $(4,5)$: $\textit{dp}[b]=5$ for $b\ge4$.
  \item After $(5,8)$: $\textit{dp}[9]=\max(5,\textit{dp}[4]+8)=13$
        (buy the $4$ and the $5$).
  \item Final $\textit{dp}[10]=13$ -- books priced $4$ and $5$ give
        $5+8=13$ pages within budget $10$.
\end{itemize}

\subsection*{C++ implementation}

\begin{lstlisting}
#include <bits/stdc++.h>
using namespace std;

int main() {
    int n, x;
    cin >> n >> x;
    vector<int> price(n), pages(n);
    for (int& p : price) cin >> p;
    for (int& s : pages) cin >> s;

    vector<int> dp(x + 1, 0);
    for (int i = 0; i < n; i++)
        for (int b = x; b >= price[i]; b--)        // downward: 0/1
            dp[b] = max(dp[b], dp[b - price[i]] + pages[i]);

    cout << dp[x] << "\n";
    return 0;
}
\end{lstlisting}

\begin{complexitybox}
\textbf{Time Complexity.} $O(nx)$ -- each book sweeps the budget once.
\textbf{Space Complexity.} $O(x)$ with the rolling 1D array.
\end{complexitybox}

\begin{takeawaybox}
Book Shop is 0/1 knapsack in disguise; the only thing to get right is the
downward inner loop that forbids reusing an item. Contrast with the coin
problems' upward loop (items reusable). Loop direction is the recurring
knapsack tell.
\end{takeawaybox}

\section{Array Description}
\label{sec:cses-array-description}

\problemstatement{An array of $n$ integers, each between $1$ and $m$, is
partially known: some positions hold a value, others are $0$ (unknown).
Count the ways to fill the unknowns so that adjacent elements differ by at
most $1$. Output modulo $10^9+7$.}

\begin{patternbox}
\emph{``Count fillings where adjacent elements are constrained''} is a DP
where the state carries the \emph{value at the current position}. Let
$\textit{dp}[i][v]$ be the number of ways to fill the prefix ending with
value $v$ at position $i$; the next position may hold $v-1,v,$ or $v+1$.
\end{patternbox}

\subsection*{State and transition}

$\textit{dp}[i][v]$ = ways to fill positions $1\ldots i$ such that
position $i$ equals $v$ and all adjacency constraints hold. From a valid
prefix ending in value $u$ at $i-1$, position $i$ may take
$u-1,u,u+1$; equivalently a value $v$ at $i$ can follow $v-1,v,v+1$ at
$i-1$:
\[
  \textit{dp}[i][v]=\textit{dp}[i-1][v-1]+\textit{dp}[i-1][v]+
  \textit{dp}[i-1][v+1].
\]
If position $i$ is fixed to some $a_i\ne0$, only $v=a_i$ is allowed (all
other $\textit{dp}[i][v]=0$); if it is $0$, every $v\in[1,m]$ is allowed.

\begin{ideabox}[Carry the Last Value in the State]
Adjacency constraints couple neighbours, so the DP state must remember the
value just placed. With that, each transition is a local three-way sum --
the $\pm1$ neighbourhood -- and fixed positions simply prune the allowed
$v$.
\end{ideabox}

\subsection*{Algorithm}

\begin{enumerate}
  \item Initialise row $1$: if $a_1$ fixed, $\textit{dp}[1][a_1]=1$;
        else $\textit{dp}[1][v]=1$ for all $v$.
  \item For $i=2\ldots n$ and each allowed $v$ at $i$:
        $\textit{dp}[i][v]=\textit{dp}[i-1][v-1]+\textit{dp}[i-1][v]+
        \textit{dp}[i-1][v+1]$.
  \item Answer $\sum_{v=1}^{m}\textit{dp}[n][v]$.
\end{enumerate}

\subsection*{Dry run}

\dryrun{$n=3$, $m=5$, array $[2,0,2]$.}

\begin{itemize}
  \item Row $1$: only $\textit{dp}[1][2]=1$.
  \item Row $2$ ($0$, free): $\textit{dp}[2][1]=\textit{dp}[1][2]=1$
        (from $2$), $\textit{dp}[2][2]=1$, $\textit{dp}[2][3]=1$; others
        $0$.
  \item Row $3$ fixed to $2$: $\textit{dp}[3][2]=\textit{dp}[2][1]+
        \textit{dp}[2][2]+\textit{dp}[2][3]=3$. Answer $3$.
\end{itemize}

\subsection*{C++ implementation}

\begin{lstlisting}
#include <bits/stdc++.h>
using namespace std;
const int MOD = 1e9 + 7;

int main() {
    int n, m;
    cin >> n >> m;
    vector<int> a(n + 1);
    for (int i = 1; i <= n; i++) cin >> a[i];

    vector<vector<long long>> dp(n + 1, vector<long long>(m + 2, 0));
    for (int v = 1; v <= m; v++)
        if (a[1] == 0 || a[1] == v) dp[1][v] = 1;

    for (int i = 2; i <= n; i++)
        for (int v = 1; v <= m; v++) {
            if (a[i] != 0 && a[i] != v) continue;      // fixed position
            dp[i][v] = (dp[i-1][v-1] + dp[i-1][v] + dp[i-1][v+1]) % MOD;
        }

    long long ans = 0;
    for (int v = 1; v <= m; v++) ans = (ans + dp[n][v]) % MOD;
    cout << ans << "\n";
    return 0;
}
\end{lstlisting}

\begin{complexitybox}
\textbf{Time Complexity.} $O(nm)$ -- $n$ positions times $m$ possible
values, each an $O(1)$ transition. \textbf{Space Complexity.} $O(nm)$,
reducible to $O(m)$ with two rolling rows.
\end{complexitybox}

\begin{takeawaybox}
When adjacent elements constrain each other, promote the ``current value''
into the DP state; the transition then only touches the local
neighbourhood the constraint allows. Fixed inputs prune states rather than
complicate the recurrence.
\end{takeawaybox}

\section{Counting Towers}
\label{sec:cses-counting-towers}

\problemstatement{Build a tower of width $2$ and height $n$ out of blocks
that are $1\times1$ or $1\times2$ (blocks may be placed in either
orientation to fill the $2$-wide grid). Count the number of distinct
towers of height $n$, modulo $10^9+7$. Answer $T$ test cases.}

\begin{patternbox}
This is a row-by-row DP with a two-valued state describing the
\textbf{boundary shape} between consecutive levels. At each height the
seam between the level below and above is either \emph{split} (two
separate cells) or \emph{joined} (one $1\times2$ block spanning the
width). Transitions between these two boundary states, counted with fixed
multipliers, give the recurrence.
\end{patternbox}

\subsection*{Two boundary states}

Process the tower one unit of height at a time, tracking the shape of the
horizontal seam at the current top:
\begin{itemize}
  \item State $0$ -- the seam is \textbf{split} into two independent
        cells (the two columns are separable at this height).
  \item State $1$ -- the seam is \textbf{joined} by a single horizontal
        $1\times2$ block spanning both columns.
\end{itemize}
Counting how each state can be extended upward by one level yields:
\[
  \begin{aligned}
    \textit{dp}[i][0] &= 4\,\textit{dp}[i-1][0] + 2\,\textit{dp}[i-1][1],\\
    \textit{dp}[i][1] &= \textit{dp}[i-1][0] + \textit{dp}[i-1][1].
  \end{aligned}
\]
The constants $4$ and $2$ enumerate the ways to fill one more level given
the seam below (four fillings keep the seam split, etc.). Base:
$\textit{dp}[1][0]=\textit{dp}[1][1]=1$.

\begin{ideabox}[The State Is the Seam, Not the Whole Tower]
The only information the next level needs from everything below is the
\emph{shape of the current top seam}. Compressing the entire history into
one of two boundary states is what makes this an $O(n)$ DP rather than an
exponential enumeration -- a miniature ``broken profile'' idea, revisited
fully in Counting Tilings (Section~\ref{sec:cses-counting-tilings}).
\end{ideabox}

\subsection*{Algorithm}

\begin{enumerate}
  \item Precompute $\textit{dp}[i][0],\textit{dp}[i][1]$ for
        $i=1\ldots\max n$ using the recurrence above.
  \item For each test case $n$, answer
        $\textit{dp}[n][0]+\textit{dp}[n][1]$ modulo $10^9+7$.
\end{enumerate}

\subsection*{Dry run}

\dryrun{Small heights.}

\begin{itemize}
  \item $n=1$: $\textit{dp}[1][0]=1,\textit{dp}[1][1]=1$; total $2$
        (two cells, or one horizontal block).
  \item $n=2$: $\textit{dp}[2][0]=4\cdot1+2\cdot1=6$,
        $\textit{dp}[2][1]=1+1=2$; total $8$.
  \item $n=3$: $\textit{dp}[3][0]=4\cdot6+2\cdot2=28$,
        $\textit{dp}[3][1]=6+2=8$; total $36$.
\end{itemize}

\subsection*{C++ implementation}

\begin{lstlisting}
#include <bits/stdc++.h>
using namespace std;
const int MOD = 1e9 + 7;
const int MAXN = 1e6 + 1;

int main() {
    static long long dp0[MAXN], dp1[MAXN];
    dp0[1] = dp1[1] = 1;
    for (int i = 2; i < MAXN; i++) {
        dp0[i] = (4 * dp0[i-1] + 2 * dp1[i-1]) % MOD;
        dp1[i] = (dp0[i-1] + dp1[i-1]) % MOD;
    }

    int t;
    cin >> t;
    while (t--) {
        int n;
        cin >> n;
        cout << (dp0[n] + dp1[n]) % MOD << "\n";
    }
    return 0;
}
\end{lstlisting}

\begin{complexitybox}
\textbf{Time Complexity.} $O(\max n)$ precomputation, then $O(1)$ per
query. \textbf{Space Complexity.} $O(\max n)$ for the two DP arrays.
\end{complexitybox}

\begin{takeawaybox}
When a structure is built layer by layer, the DP state need only capture
the \emph{interface} between the built and unbuilt parts. Here two seam
states suffice; identifying that minimal boundary description is the core
skill, scaling up to bitmask profiles for wider grids.
\end{takeawaybox}

\section{Edit Distance}
\label{sec:cses-edit-distance}

\problemstatement{Given two strings, find their edit (Levenshtein)
distance: the minimum number of single-character insertions, deletions,
and substitutions needed to turn one string into the other.}

\begin{patternbox}
\emph{``Minimum operations to transform one sequence into another''} is
the canonical two-sequence grid DP. Let $\textit{dp}[i][j]$ be the edit
distance between the first $i$ characters of $A$ and the first $j$ of $B$;
the last operation is an insert, delete, or match/substitute.
\end{patternbox}

\subsection*{State and transition}

$\textit{dp}[i][j]$ = edit distance between prefixes $A[1..i]$ and
$B[1..j]$. Consider the last operation aligning the two prefixes:
\begin{itemize}
  \item If $A[i]=B[j]$, they match for free:
        $\textit{dp}[i-1][j-1]$.
  \item Otherwise substitute: $1+\textit{dp}[i-1][j-1]$.
  \item Delete $A[i]$: $1+\textit{dp}[i-1][j]$.
  \item Insert $B[j]$: $1+\textit{dp}[i][j-1]$.
\end{itemize}
\[
  \textit{dp}[i][j]=\min\!\big(\textit{dp}[i-1][j-1]+[A_i\ne B_j],\
  \textit{dp}[i-1][j]+1,\ \textit{dp}[i][j-1]+1\big).
\]
Base cases: $\textit{dp}[i][0]=i$ (delete all), $\textit{dp}[0][j]=j$
(insert all).

\begin{ideabox}[Three Operations, Three Predecessors]
Each cell chooses the cheapest of the three edits that could have produced
it, so the recurrence has exactly three incoming states. The diagonal move
is free on a character match -- that is where the strings' similarity saves
operations.
\end{ideabox}

\subsection*{Algorithm}

\begin{enumerate}
  \item Fill the first row and column with $0,1,2,\ldots$ (transform to
        or from the empty string).
  \item For $i=1\ldots|A|$, $j=1\ldots|B|$: apply the recurrence above.
  \item Answer $\textit{dp}[|A|][|B|]$.
\end{enumerate}

\subsection*{Dry run}

\dryrun{$A=$``LOVE'', $B=$``MOVIE''.}

\begin{itemize}
  \item Align: L$\to$M (substitute), O=O, V=V, insert I, E=E.
  \item The DP settles $\textit{dp}[4][5]=2$ -- one substitution
        (L$\to$M) and one insertion (I).
\end{itemize}

\subsection*{C++ implementation}

\begin{lstlisting}
#include <bits/stdc++.h>
using namespace std;

int main() {
    string a, b;
    cin >> a >> b;
    int n = a.size(), m = b.size();

    vector<vector<int>> dp(n + 1, vector<int>(m + 1, 0));
    for (int i = 0; i <= n; i++) dp[i][0] = i;
    for (int j = 0; j <= m; j++) dp[0][j] = j;

    for (int i = 1; i <= n; i++)
        for (int j = 1; j <= m; j++) {
            int cost = (a[i - 1] == b[j - 1]) ? 0 : 1;
            dp[i][j] = min({ dp[i-1][j-1] + cost,   // match/substitute
                             dp[i-1][j] + 1,        // delete
                             dp[i][j-1] + 1 });     // insert
        }
    cout << dp[n][m] << "\n";
    return 0;
}
\end{lstlisting}

\begin{complexitybox}
\textbf{Time Complexity.} $O(nm)$. \textbf{Space Complexity.} $O(nm)$,
reducible to $O(\min(n,m))$ with two rolling rows.
\end{complexitybox}

\begin{takeawaybox}
Edit Distance is the two-string grid DP every alignment problem descends
from: three edit operations map to three predecessor cells, and a matching
character makes the diagonal free. The same $(i,j)$-prefix state solves
LCS, sequence alignment, and diff.
\end{takeawaybox}

\section{Rectangle Cutting}
\label{sec:cses-rectangle-cutting}

\problemstatement{Given an $a\times b$ rectangle, repeatedly cut it with
straight horizontal or vertical cuts (each cut splits a piece into two,
along integer coordinates) until only squares remain. Find the minimum
number of cuts needed.}

\begin{patternbox}
\emph{``Minimum cuts to reduce a rectangle to squares''} is an interval
DP in two dimensions. Let $\textit{dp}[h][w]$ be the fewest cuts for an
$h\times w$ rectangle; a first cut splits it into two smaller rectangles
whose costs are already known.
\end{patternbox}

\subsection*{State and transition}

$\textit{dp}[h][w]$ = minimum cuts for an $h\times w$ piece. If $h=w$ it
is already a square: $\textit{dp}[h][h]=0$. Otherwise the first cut is
either horizontal (splitting height into $i$ and $h-i$) or vertical
(splitting width into $j$ and $w-j$), each costing one plus the two
sub-costs:
\[
  \textit{dp}[h][w]=1+\min\!\Big(
  \min_{1\le i<h}\!\big(\textit{dp}[i][w]+\textit{dp}[h-i][w]\big),\
  \min_{1\le j<w}\!\big(\textit{dp}[h][j]+\textit{dp}[h][w-j]\big)\Big).
\]

\begin{ideabox}[Try Every First Cut]
As in all interval DP, we do not guess the optimal cut -- we try every
possible first cut and take the minimum, trusting that sub-rectangles are
solved optimally. Squares terminate the recursion at zero cost.
\end{ideabox}

\subsection*{Algorithm}

\begin{enumerate}
  \item For all $s$, $\textit{dp}[s][s]=0$.
  \item For increasing $h,w$ with $h\ne w$: take the best over all
        horizontal splits $i$ and vertical splits $j$, add $1$.
  \item Answer $\textit{dp}[a][b]$.
\end{enumerate}

\subsection*{Dry run}

\dryrun{$a=3$, $b=5$.}

\begin{itemize}
  \item $\textit{dp}[3][3]=0$. $\textit{dp}[2][3]=1+\min$ over cuts
        $=1+\textit{dp}[2][2]+\textit{dp}[2][1]$-type options $=2$ (cut
        off a $2\times2$, then split the $2\times1$).
  \item Working up, $\textit{dp}[3][5]=3$ -- e.g.\ cut a $3\times5$ into
        $3\times3$ and $3\times2$, then the $3\times2$ into a $2\times2$
        and pieces, totalling three cuts.
\end{itemize}

\subsection*{C++ implementation}

\begin{lstlisting}
#include <bits/stdc++.h>
using namespace std;

int main() {
    int a, b;
    cin >> a >> b;
    vector<vector<int>> dp(a + 1, vector<int>(b + 1, 0));

    for (int h = 1; h <= a; h++)
        for (int w = 1; w <= b; w++) {
            if (h == w) { dp[h][w] = 0; continue; }
            int best = INT_MAX;
            for (int i = 1; i < h; i++)                 // horizontal cut
                best = min(best, dp[i][w] + dp[h - i][w]);
            for (int j = 1; j < w; j++)                 // vertical cut
                best = min(best, dp[h][j] + dp[h][w - j]);
            dp[h][w] = best + 1;
        }
    cout << dp[a][b] << "\n";
    return 0;
}
\end{lstlisting}

\begin{complexitybox}
\textbf{Time Complexity.} $O(ab\,(a+b))$ -- each of the $ab$ states tries
$O(a+b)$ cut positions. \textbf{Space Complexity.} $O(ab)$.
\end{complexitybox}

\begin{takeawaybox}
Interval DP in 2D: enumerate every first cut, recurse on the two pieces,
take the minimum. The ``try all split points'' pattern recurs in Removal
Game and any problem where a structure is optimally divided into
independent parts.
\end{takeawaybox}

\section{Money Sums}
\label{sec:cses-money-sums}

\problemstatement{You have $n$ coins with given values. Determine all the
distinct sums you can form using \emph{some subset} of the coins (each
coin used at most once). Output how many distinct sums are possible and
list them.}

\begin{patternbox}
\emph{``Which totals are achievable from a subset''} is the boolean
\textbf{subset-sum} DP. Let $\textit{dp}[s]$ be true if some subset sums
to exactly $s$; each coin either joins a subset or not, a $0/1$ choice.
\end{patternbox}

\subsection*{State and transition}

$\textit{dp}[s]$ = whether sum $s$ is reachable using a subset of the
coins seen so far, with $\textit{dp}[0]=$ true (the empty subset). Adding
a coin of value $c$ makes any previously reachable $s$ also reach $s+c$.
Using a 1D array we iterate $s$ \textbf{downward} (0/1 knapsack rule) so a
coin is not reused:
\[
  \textit{dp}[s]\gets \textit{dp}[s]\ \lor\ \textit{dp}[s-c].
\]

\begin{ideabox}[Reachability, Not Counting]
Money Sums only asks \emph{whether} a sum is attainable, so the DP stores
booleans and combines with OR. The total possible sum is
$\sum c_i\le10^5$, so a single boolean array over $[0,\text{total}]$
suffices, and the answer is every index set to true.
\end{ideabox}

\subsection*{Algorithm}

\begin{enumerate}
  \item Let $S=\sum c_i$. Set $\textit{dp}[0]=$ true.
  \item For each coin $c$: for $s=S$ down to $c$,
        $\textit{dp}[s]\mathrel{|}=\textit{dp}[s-c]$.
  \item Collect all $s\ge1$ with $\textit{dp}[s]$ true; their count and
        list are the answer.
\end{enumerate}

\subsection*{Dry run}

\dryrun{Coins $\{4,2,2\}$; total $S=8$.}

\begin{itemize}
  \item Start $\textit{dp}[0]=$T. After coin $4$: $\{0,4\}$.
  \item After a $2$: $\{0,2,4,6\}$. After the second $2$:
        $\{0,2,4,6,8\}$.
  \item Non-zero reachable sums: $2,4,6,8$ -- four distinct sums.
\end{itemize}

\subsection*{C++ implementation}

\begin{lstlisting}
#include <bits/stdc++.h>
using namespace std;

int main() {
    int n;
    cin >> n;
    vector<int> coins(n);
    int total = 0;
    for (int& c : coins) { cin >> c; total += c; }

    vector<char> dp(total + 1, 0);
    dp[0] = 1;
    for (int c : coins)
        for (int s = total; s >= c; s--)           // downward: 0/1
            if (dp[s - c]) dp[s] = 1;

    vector<int> sums;
    for (int s = 1; s <= total; s++) if (dp[s]) sums.push_back(s);

    cout << sums.size() << "\n";
    for (int s : sums) cout << s << " ";
    cout << "\n";
    return 0;
}
\end{lstlisting}

\begin{complexitybox}
\textbf{Time Complexity.} $O(n\cdot S)$ where $S=\sum c_i$.
\textbf{Space Complexity.} $O(S)$ for the boolean array.
\end{complexitybox}

\begin{takeawaybox}
Subset-sum reachability is 0/1 knapsack with booleans and OR. The downward
loop again enforces ``each coin once.'' This reachability array is the
engine behind Two Sets II (next), where we \emph{count} the subsets rather
than merely test existence.
\end{takeawaybox}

\section{Removal Game}
\label{sec:cses-removal-game}

\problemstatement{An array of $n$ numbers is on a table. Two players
alternate; on each turn a player removes either the leftmost or rightmost
element and adds its value to their score. Both play optimally to maximise
their own score. Find the first player's final score.}

\begin{patternbox}
\emph{``Two players take from the ends, both optimal''} is an interval
game DP. Let $\textit{dp}[l][r]$ be the best score \emph{difference} (current
player minus opponent) achievable on the subarray $[l,r]$; the current
player takes an end and then \emph{becomes} the opponent on the rest.
\end{patternbox}

\subsection*{State and transition via score difference}

Track $\textit{dp}[l][r]$ = the maximum of (current player's score $-$
opponent's score) over the subarray $a[l..r]$, assuming optimal play. The
current player takes $a[l]$ or $a[r]$; after that the opponent plays
optimally on the remaining interval, so their advantage is subtracted:
\[
  \textit{dp}[l][r]=\max\!\big(a[l]-\textit{dp}[l+1][r],\
  a[r]-\textit{dp}[l][r-1]\big),\qquad \textit{dp}[l][l]=a[l].
\]
If $\textit{total}=\sum a_i$ and $d=\textit{dp}[0][n-1]$ is the final
difference, the first player's score is $(\textit{total}+d)/2$.

\begin{ideabox}[Score Difference Turns Two Objectives into One]
Modelling the game by the \emph{difference} folds ``maximise mine'' and
``opponent minimises mine'' into a single recurrence: whatever the
opponent gains counts against me, so I subtract their optimal advantage on
the remainder. This zero-sum trick avoids storing two separate scores.
\end{ideabox}

\subsection*{Algorithm}

\begin{enumerate}
  \item $\textit{dp}[i][i]=a[i]$.
  \item For increasing interval length, $\textit{dp}[l][r]=\max(a[l]-
        \textit{dp}[l+1][r],\ a[r]-\textit{dp}[l][r-1])$.
  \item First player's score $=(\textit{total}+\textit{dp}[0][n-1])/2$.
\end{enumerate}

\subsection*{Dry run}

\dryrun{$a=[4,5,1,3]$, total $=13$.}

\begin{itemize}
  \item Length 1: $\textit{dp}[i][i]=a[i]$.
  \item Length 2: $\textit{dp}[0][1]=\max(4-5,5-4)=1$;
        $\textit{dp}[2][3]=\max(1-3,3-1)=2$; $\textit{dp}[1][2]=\max(5-1,
        1-5)=4$.
  \item Full: $\textit{dp}[0][3]=\max(4-\textit{dp}[1][3],
        3-\textit{dp}[0][2])$. With $\textit{dp}[1][3]=\max(5-2,3-4)=3$
        and $\textit{dp}[0][2]=\max(4-4,1-1)=0$: $\max(4-3,3-0)=3$.
  \item First player's score $=(13+3)/2=8$.
\end{itemize}

\subsection*{C++ implementation}

\begin{lstlisting}
#include <bits/stdc++.h>
using namespace std;

int main() {
    int n;
    cin >> n;
    vector<long long> a(n);
    long long total = 0;
    for (auto& x : a) { cin >> x; total += x; }

    vector<vector<long long>> dp(n, vector<long long>(n, 0));
    for (int i = 0; i < n; i++) dp[i][i] = a[i];

    for (int len = 2; len <= n; len++)
        for (int l = 0; l + len - 1 < n; l++) {
            int r = l + len - 1;
            dp[l][r] = max(a[l] - dp[l + 1][r], a[r] - dp[l][r - 1]);
        }

    cout << (total + dp[0][n - 1]) / 2 << "\n";
    return 0;
}
\end{lstlisting}

\begin{complexitybox}
\textbf{Time Complexity.} $O(n^2)$ -- one $O(1)$ transition per interval.
\textbf{Space Complexity.} $O(n^2)$ for the interval table.
\end{complexitybox}

\begin{takeawaybox}
For symmetric two-player, zero-sum games, track the \emph{score
difference} from the mover's perspective and subtract the opponent's
optimal remainder. This collapses adversarial optimisation into a single
maximisation -- the standard move for end-taking and Nim-like game DPs.
\end{takeawaybox}

\section{Two Sets II}
\label{sec:cses-two-sets-ii}

\problemstatement{Count the number of ways to partition the numbers
$\{1,2,\ldots,n\}$ into two subsets with \emph{equal} sums. Two partitions
that are reflections of each other count as the same. Output modulo
$10^9+7$.}

\begin{patternbox}
\emph{``Split $\{1..n\}$ into two equal-sum halves, count the ways''} is a
\textbf{counting subset-sum}. The total is
$T=\frac{n(n+1)}{2}$; if $T$ is odd the answer is $0$, else we count
subsets summing to $T/2$, then divide by $2$ to remove the mirror
duplication.
\end{patternbox}

\subsection*{State and transition}

Let the target be $t=T/2$. Define $\textit{dp}[s]$ = number of subsets of
$\{1,\ldots,n\}$ (processed one number at a time) that sum to $s$, with
$\textit{dp}[0]=1$. Each number $k$ is either in the subset or not -- a
$0/1$ choice -- so with the downward inner loop:
\[
  \textit{dp}[s]\mathrel{+}=\textit{dp}[s-k]\quad\text{for } s=t\ \text{down to}\ k.
\]
Each valid partition is counted twice (subset $A$ and its complement
$B$), so the final answer is $\textit{dp}[t]/2$ modulo $10^9+7$ (using the
modular inverse of $2$).

\begin{ideabox}[Count Subsets, Then Halve for the Mirror]
Counting subsets that reach exactly half the total counts every partition
from both sides. Dividing by $2$ (via the modular inverse, since we are
mod a prime) removes that symmetry. The odd-total shortcut avoids
pointless work when no equal split exists.
\end{ideabox}

\subsection*{Algorithm}

\begin{enumerate}
  \item Compute $T=n(n+1)/2$; if $T$ is odd, output $0$.
  \item Set $t=T/2$, $\textit{dp}[0]=1$. For $k=1\ldots n-1$ (one number
        may be fixed to a side to pre-halve, or process all and divide):
        update $\textit{dp}[s]\mathrel{+}=\textit{dp}[s-k]$ downward.
  \item Output $\textit{dp}[t]$, halved via the modular inverse of $2$.
\end{enumerate}

(Alternatively, process $k=1\ldots n-1$ only -- fixing $n$ to one side --
which directly avoids the double count without a division.)

\subsection*{Dry run}

\dryrun{$n=7$, $T=28$, target $t=14$.}

\begin{itemize}
  \item Subsets of $\{1..7\}$ summing to $14$ include $\{7,6,1\},
        \{7,5,2\},\{7,4,3\},\{7,4,2,1\},\ldots$
  \item The counting DP gives $\textit{dp}[14]=8$; halving yields the
        answer $4$ distinct partitions.
\end{itemize}

\subsection*{C++ implementation}

\begin{lstlisting}
#include <bits/stdc++.h>
using namespace std;
const long long MOD = 1e9 + 7;

int main() {
    int n;
    cin >> n;
    long long total = (long long)n * (n + 1) / 2;
    if (total % 2 != 0) { cout << 0 << "\n"; return 0; }

    int t = total / 2;
    vector<long long> dp(t + 1, 0);
    dp[0] = 1;
    // fix element n on one side: process 1..n-1, no division needed
    for (int k = 1; k <= n - 1; k++)
        for (int s = t; s >= k; s--)
            dp[s] = (dp[s] + dp[s - k]) % MOD;

    cout << dp[t] << "\n";
    return 0;
}
\end{lstlisting}

\begin{complexitybox}
\textbf{Time Complexity.} $O(n\cdot T)=O(n^3)$ in the worst sense, but
$T=O(n^2)$ so it is $O(n\cdot n^2)$; practically $O(n\cdot t)$ with
$t\approx n^2/4$. \textbf{Space Complexity.} $O(t)=O(n^2)$.
\end{complexitybox}

\begin{takeawaybox}
Two Sets II is counting subset-sum with a symmetry fix. The clean way to
avoid the double count is to fix one fixed element (here $n$) to a single
side and only distribute the rest -- turning ``divide by two'' into ``count
directly.'' Recognise equal-partition problems as subset-sum to half the
total.
\end{takeawaybox}

\section{Increasing Subsequence}
\label{sec:cses-increasing-subsequence}

\problemstatement{Given an array of $n$ integers, find the length of the
longest strictly increasing subsequence (LIS): the longest sequence of
indices $i_1<i_2<\cdots$ with strictly increasing values.}

\begin{patternbox}
\emph{``Longest increasing subsequence''} has an elegant
$O(n\log n)$ solution built on \textbf{patience sorting}: maintain the
smallest possible tail value for an increasing subsequence of each
length, and binary-search where each new element fits.
\end{patternbox}

\subsection*{The tails array}

Keep an array \textit{tails}, where $\textit{tails}[k]$ is the smallest
value that can end an increasing subsequence of length $k+1$ seen so far.
This array is always sorted. For each new value $x$, find the first tail
$\ge x$ (a \texttt{lower\_bound}) and overwrite it with $x$: this either
extends the longest subsequence (if $x$ is larger than every tail,
append) or improves a length-$k$ subsequence's ending to a smaller value,
leaving room for future growth. The LIS length is the final size of
\textit{tails}.

\begin{ideabox}[Keep the Smallest Tail per Length]
Replacing the first tail $\ge x$ with $x$ never shortens any achievable
length and keeps every tail as small as possible, maximising the chance of
future extension. The binary search is what turns the naive $O(n^2)$ DP
into $O(n\log n)$.
\end{ideabox}

\subsection*{Algorithm}

\begin{enumerate}
  \item Start with an empty \textit{tails}.
  \item For each $x$: binary-search the first index with
        $\textit{tails}[i]\ge x$. If none, append $x$ (LIS grew); else
        set $\textit{tails}[i]=x$.
  \item Answer is the size of \textit{tails}.
\end{enumerate}

\subsection*{Dry run}

\dryrun{$a=[7,3,5,3,6,2,9,8]$.}

\begin{itemize}
  \item $7\to[7]$; $3$ replaces $7\to[3]$; $5$ appends $\to[3,5]$; $3$
        replaces $5$? No -- lower\_bound of $3$ is index $0$ ($3\ge3$),
        replace $\to[3,5]$ (unchanged in effect).
  \item $6$ appends $\to[3,5,6]$; $2$ replaces $3\to[2,5,6]$; $9$
        appends $\to[2,5,6,9]$; $8$ replaces $9\to[2,5,6,8]$.
  \item LIS length $=4$ (e.g.\ $3,5,6,9$).
\end{itemize}

\subsection*{C++ implementation}

\begin{lstlisting}
#include <bits/stdc++.h>
using namespace std;

int main() {
    int n;
    cin >> n;
    vector<int> a(n);
    for (int& x : a) cin >> x;

    vector<int> tails;
    for (int x : a) {
        auto it = lower_bound(tails.begin(), tails.end(), x);  // first >= x
        if (it == tails.end()) tails.push_back(x);   // extend LIS
        else *it = x;                                // shrink a tail
    }
    cout << tails.size() << "\n";
    return 0;
}
\end{lstlisting}

\begin{complexitybox}
\textbf{Time Complexity.} $O(n\log n)$ -- one binary search per element.
\textbf{Space Complexity.} $O(n)$ for the tails array. (A simpler
$O(n^2)$ DP with $\textit{dp}[i]=1+\max_{j<i,a_j<a_i}\textit{dp}[j]$ also
solves it for small $n$.)
\end{complexitybox}

\begin{takeawaybox}
LIS in $O(n\log n)$ is the patience-sorting tails trick: keep the minimal
ending value for each length, binary-search each insertion. For strictly
increasing use \texttt{lower\_bound}; for non-decreasing use
\texttt{upper\_bound} -- a one-word change that trips up many.
\end{takeawaybox}

\section{Projects}
\label{sec:cses-projects}

\problemstatement{There are $n$ projects, each with a start day, an end
day, and a reward. Attending a project occupies all days from its start to
its end inclusive, so two attended projects may not overlap. Maximise the
total reward.}

\begin{patternbox}
\emph{``Non-overlapping intervals, maximise total value''} is
\textbf{weighted interval scheduling}. Sort projects by end day; for each,
either skip it or take it and jump to the best solution that finishes
before its start -- found by binary search.
\end{patternbox}

\subsection*{State and transition}

Sort the projects by end day. Let $\textit{dp}[i]$ be the maximum reward
using only the first $i$ projects (in sorted order). For project $i$ with
start $s_i$ and reward $p_i$, either skip it ($\textit{dp}[i-1]$) or take
it: add $p_i$ to the best reward achievable among projects that end
\emph{strictly before} $s_i$. Let $j$ be the number of projects whose end
day $<s_i$ (found by binary search on the sorted end days):
\[
  \textit{dp}[i]=\max\big(\textit{dp}[i-1],\ p_i+\textit{dp}[j]\big).
\]

\begin{ideabox}[Sort by End, Binary-Search the Predecessor]
Sorting by end day means that when we consider a project, every compatible
earlier project has already been folded into some $\textit{dp}[j]$. Binary
search finds the latest non-conflicting project in $O(\log n)$, turning an
$O(n^2)$ scan into $O(n\log n)$.
\end{ideabox}

\subsection*{Algorithm}

\begin{enumerate}
  \item Sort projects by end day; let $\textit{ends}[]$ be their end
        days.
  \item $\textit{dp}[0]=0$. For $i=1\ldots n$: binary-search $j$ = number
        of projects with end $<s_i$; set
        $\textit{dp}[i]=\max(\textit{dp}[i-1], p_i+\textit{dp}[j])$.
  \item Answer $\textit{dp}[n]$.
\end{enumerate}

\subsection*{Dry run}

\dryrun{Projects (start, end, reward): $(2,4,7),(1,2,4),(5,6,3)$.}

\begin{itemize}
  \item Sorted by end: $(1,2,4),(2,4,7),(5,6,3)$.
  \item $\textit{dp}[1]=4$. Project $(2,4,7)$: projects ending $<2$ is
        the first (end $2$? no, need $<2$, so $j=0$):
        $\textit{dp}[2]=\max(4,7+0)=7$.
  \item Project $(5,6,3)$: projects ending $<5$ are both, $j=2$:
        $\textit{dp}[3]=\max(7,3+7)=10$. Answer $10$.
\end{itemize}

\subsection*{C++ implementation}

\begin{lstlisting}
#include <bits/stdc++.h>
using namespace std;

int main() {
    int n;
    cin >> n;
    vector<array<long long,3>> p(n);           // {end, start, reward}
    for (int i = 0; i < n; i++) {
        long long a, b, x;
        cin >> a >> b >> x;
        p[i] = {b, a, x};
    }
    sort(p.begin(), p.end());                  // by end day

    vector<long long> ends(n);
    for (int i = 0; i < n; i++) ends[i] = p[i][0];

    vector<long long> dp(n + 1, 0);
    for (int i = 1; i <= n; i++) {
        long long start = p[i-1][1], reward = p[i-1][2];
        // number of projects with end < start
        int j = lower_bound(ends.begin(), ends.end(), start) - ends.begin();
        dp[i] = max(dp[i-1], reward + dp[j]);
    }
    cout << dp[n] << "\n";
    return 0;
}
\end{lstlisting}

\begin{complexitybox}
\textbf{Time Complexity.} $O(n\log n)$ -- sorting plus a binary search per
project. \textbf{Space Complexity.} $O(n)$.
\end{complexitybox}

\begin{takeawaybox}
Weighted interval scheduling = sort by end, then ``skip or take with a
binary-searched predecessor.'' It generalises the unweighted greedy
(which just maximises count) to rewards, where greed fails and DP is
required.
\end{takeawaybox}

\section{Elevator Rides}
\label{sec:cses-elevator-rides}

\problemstatement{There are $n$ people with given weights, and an elevator
with maximum weight $x$. Find the minimum number of elevator rides needed
to carry everyone (people may be grouped arbitrarily, but each ride's
total weight must not exceed $x$).}

\begin{patternbox}
With $n\le20$, ``partition people into fewest groups under a capacity''
screams \textbf{bitmask DP}: the state is the \emph{subset} of people
already transported. The clever part is what to store -- not just the ride
count, but how full the last ride is.
\end{patternbox}

\subsection*{A two-part state per subset}

For each subset \textit{mask} of people, store a pair
$\textit{dp}[\textit{mask}]=(\textit{rides},\textit{lastWeight})$: the
minimum number of rides to carry exactly that subset, and among solutions
with that few rides, the smallest weight currently loaded in the last
(still-open) ride. Compare pairs lexicographically -- fewer rides first,
then a lighter last ride (which leaves the most room). Transition by adding
one new person $p$ to $\textit{mask}$:
\begin{itemize}
  \item If $p$ fits in the current last ride
        ($\textit{lastWeight}+w_p\le x$): same ride count, last weight
        grows by $w_p$.
  \item Otherwise start a new ride: rides $+1$, last weight $=w_p$.
\end{itemize}

\begin{ideabox}[Store Rides \emph{and} the Last Ride's Load]
Minimising rides alone is not enough, because two subsets with equal ride
counts differ in how much room remains -- and that room affects future
additions. Pairing the ride count with the least-loaded last ride makes
the greedy-within-DP choice optimal.
\end{ideabox}

\subsection*{Algorithm}

\begin{enumerate}
  \item $\textit{dp}[0]=(1,0)$ -- one empty ride, zero load.
  \item For each mask and each person $p\notin\textit{mask}$: compute
        the candidate pair for $\textit{mask}\cup\{p\}$ using the rule
        above; keep the lexicographically smaller.
  \item Answer is the ride count of $\textit{dp}[\text{full mask}]$.
\end{enumerate}

\subsection*{Dry run}

\dryrun{Weights $[2,3,3]$, capacity $x=5$.}

\begin{itemize}
  \item $\textit{dp}[\{0\}]=(1,2)$, $\textit{dp}[\{1\}]=(1,3)$.
  \item Add person $2$ ($w=3$) to $\{0\}$ (last $2$): $2+3=5\le5$, same
        ride: $(1,5)$. To $\{1\}$ (last $3$): $3+3=6>5$, new ride:
        $(2,3)$.
  \item $\textit{dp}[\{0,2\}]=(1,5)$, then adding person $1$
        ($3$): $5+3>5$, new ride $(2,3)$. Full-set answer: $2$ rides
        (e.g.\ $\{2,3\}$ and $\{3\}$).
\end{itemize}

\subsection*{C++ implementation}

\begin{lstlisting}
#include <bits/stdc++.h>
using namespace std;

int main() {
    int n, x;
    cin >> n >> x;
    vector<int> w(n);
    for (int& v : w) cin >> v;

    int full = 1 << n;
    // dp[mask] = {rides, weight in last ride}
    vector<pair<int,int>> dp(full, {n + 1, 0});
    dp[0] = {1, 0};

    for (int mask = 1; mask < full; mask++) {
        for (int p = 0; p < n; p++) {
            if (!(mask & (1 << p))) continue;
            auto prev = dp[mask ^ (1 << p)];       // state without p
            if (prev.second + w[p] <= x)           // fits in last ride
                prev.second += w[p];
            else                                   // needs a new ride
                prev = {prev.first + 1, w[p]};
            dp[mask] = min(dp[mask], prev);        // lexicographic
        }
    }
    cout << dp[full - 1].first << "\n";
    return 0;
}
\end{lstlisting}

\begin{complexitybox}
\textbf{Time Complexity.} $O(2^n\cdot n)$ -- each of $2^n$ subsets tries
adding each of $n$ people. \textbf{Space Complexity.} $O(2^n)$ for the DP
table.
\end{complexitybox}

\begin{takeawaybox}
Bitmask DP fits partition problems with $n\le20$. The lesson beyond the
mask is the \emph{compound state}: pairing ``rides so far'' with ``room
left in the current ride'' captures exactly the extra information a greedy
packing needs -- a recurring trick when a count alone loses too much.
\end{takeawaybox}

\section{Counting Tilings}
\label{sec:cses-counting-tilings}

\problemstatement{Count the number of ways to tile an $n\times m$ grid
completely with $1\times2$ dominoes (each placed horizontally or
vertically), modulo $10^9+7$.}

\begin{patternbox}
\emph{``Domino tilings of a grid''} with one small dimension ($m\le10$) is
solved by \textbf{broken-profile (bitmask) DP}: sweep column by column,
and let the state be the bitmask of cells in the current column that are
already occupied by a horizontal domino protruding from the previous
column.
\end{patternbox}

\subsection*{The profile between columns}

Fill the grid one column at a time. When moving from column $c$ to column
$c+1$, the only thing the next column needs to know is which of its cells
are \emph{already filled} by horizontal dominoes sticking out of column
$c$. That set of $n$ bits is the \textbf{profile}. For each incoming
profile \textit{mask}, we enumerate every legal way to complete column $c$
using vertical dominoes (within the column) and horizontal dominoes
(extending into column $c+1$), producing an outgoing profile.

\begin{ideabox}[Recursively Fill One Column]
Process the current column cell by cell from top to bottom. A cell already
covered by the incoming profile is skipped. An empty cell must be covered
either by a vertical domino with the cell below it, or by a horizontal
domino that sets the corresponding bit in the \emph{next} column's
profile. Enumerating these choices generates all valid transitions.
\end{ideabox}

\subsection*{Algorithm}

\begin{enumerate}
  \item Ensure $n$ (the profile height) is the smaller dimension; if not,
        swap $n$ and $m$.
  \item Start with $\textit{cur}[0]=1$ (empty profile before column $0$).
  \item For each column, for each profile with nonzero count, recursively
        fill the column to produce all next-column profiles, accumulating
        counts.
  \item After the last column, the answer is $\textit{cur}[0]$ (nothing
        protrudes past the grid).
\end{enumerate}

\subsection*{Dry run}

\dryrun{$2\times2$ grid.}

\begin{itemize}
  \item Column $0$, incoming profile $00$: either two horizontals (both
        cells extend right, next profile $11$) or one vertical (next
        profile $00$).
  \item Column $1$: profile $11$ means both cells already filled -- valid,
        next $00$. Profile $00$ needs a vertical -- valid, next $00$.
  \item Total tilings of $2\times2$: $2$ (two horizontals or two
        verticals).
\end{itemize}

\subsection*{C++ implementation}

\begin{lstlisting}
#include <bits/stdc++.h>
using namespace std;
const long long MOD = 1e9 + 7;

int n, m;
vector<long long> cur, nxt;

// fill column: row = current row, mask = incoming profile,
// newmask = outgoing profile being built, ways = count to propagate
void fill(int row, int mask, int newmask, long long ways) {
    if (row == n) { nxt[newmask] = (nxt[newmask] + ways) % MOD; return; }
    if (mask & (1 << row)) {                     // already filled
        fill(row + 1, mask, newmask, ways);
        return;
    }
    // place a horizontal domino: extends into next column
    fill(row + 1, mask, newmask | (1 << row), ways);
    // place a vertical domino covering row and row+1
    if (row + 1 < n && !(mask & (1 << (row + 1))))
        fill(row + 2, mask, newmask, ways);
}

int main() {
    cin >> n >> m;
    if (n > m) swap(n, m);                        // keep profile small
    int full = 1 << n;
    cur.assign(full, 0);
    cur[0] = 1;
    for (int col = 0; col < m; col++) {
        nxt.assign(full, 0);
        for (int mask = 0; mask < full; mask++)
            if (cur[mask]) fill(0, mask, 0, cur[mask]);
        cur = nxt;
    }
    cout << cur[0] << "\n";                       // nothing protrudes out
    return 0;
}
\end{lstlisting}

\begin{complexitybox}
\textbf{Time Complexity.} $O(m\cdot 2^n\cdot n)$ roughly -- each column
processes $2^n$ profiles, each filled in $O(n)$ recursive steps.
\textbf{Space Complexity.} $O(2^n)$ for the two profile arrays.
\end{complexitybox}

\begin{takeawaybox}
Broken-profile DP tiles a grid by carrying, as state, only the boundary
between the filled and unfilled regions -- one column's protrusion mask.
It is Counting Towers (Section~\ref{sec:cses-counting-towers}) generalised
from two boundary states to $2^n$, and the standard tool for domino and
polyomino tiling counts.
\end{takeawaybox}

\section{Counting Numbers}
\label{sec:cses-counting-numbers}

\problemstatement{Count the integers in the range $[a,b]$ that do not
contain two adjacent equal digits (e.g.\ $33$ and $121\underline{22}$ are
excluded). Here $0\le a\le b\le10^{18}$.}

\begin{patternbox}
\emph{``Count numbers up to $N$ with a digit property''} is
\textbf{digit DP}. Reduce $[a,b]$ to $f(b)-f(a-1)$ where $f(x)$ counts
valid numbers in $[0,x]$, then build $x$'s digits left to right, tracking
just enough state to enforce ``no two adjacent equal digits.''
\end{patternbox}

\subsection*{State for digit DP}

Write $x$ as a digit string and fill positions left to right. The state
is:
\begin{itemize}
  \item \textbf{position} in the digit string,
  \item \textbf{previous digit} placed (to forbid an equal neighbour),
  \item \textbf{tight}: whether the prefix so far equals $x$'s prefix (if
        so, the next digit is capped by $x$'s digit),
  \item \textbf{started}: whether we have placed a nonzero digit yet (to
        handle leading zeros, which must not count as ``adjacent equal'').
\end{itemize}
At each position we try every allowed next digit $d$ (up to the cap if
tight), skipping $d$ equal to the previous digit once the number has
started. Memoise over (position, previous, started) for the non-tight
states.

\begin{ideabox}[Tight and Started Are the Two Universal Digit-DP Flags]
\textbf{Tight} keeps the enumeration within $[0,x]$; it is true only along
the single prefix that hugs $x$. \textbf{Started} distinguishes real
leading digits from padding zeros, so a leading zero does not falsely
collide with the next digit. Almost every digit-DP carries these two.
\end{ideabox}

\subsection*{Algorithm}

\begin{enumerate}
  \item Define $f(x)$ = count of valid numbers in $[0,x]$ by digit DP
        over $x$'s digits with state (pos, prev, tight, started).
  \item At each position try digits $0\ldots(\text{cap})$; skip $d=$
        prev when started; recurse with updated flags.
  \item Answer $f(b)-f(a-1)$ (with $f(-1)=0$ when $a=0$).
\end{enumerate}

\subsection*{Dry run}

\dryrun{Count valid numbers in $[0,23]$.}

\begin{itemize}
  \item Invalid two-digit numbers $\le23$ with equal adjacent digits:
        $11,22$.
  \item All $24$ numbers $0\ldots23$ minus the $2$ invalid ones $=22$
        valid numbers.
  \item The digit DP arrives at the same $22$ by never placing a digit
        equal to its left neighbour.
\end{itemize}

\subsection*{C++ implementation}

\begin{lstlisting}
#include <bits/stdc++.h>
using namespace std;

string s;
long long memo[20][11][2];      // [pos][prev+1][started], for tight=false

long long solve(int pos, int prev, bool tight, bool started) {
    if (pos == (int)s.size()) return 1;          // one valid number formed
    if (!tight && memo[pos][prev + 1][started] != -1)
        return memo[pos][prev + 1][started];

    int cap = tight ? (s[pos] - '0') : 9;
    long long total = 0;
    for (int d = 0; d <= cap; d++) {
        if (started && d == prev) continue;      // no equal adjacent digits
        bool nstarted = started || (d > 0);
        int nprev = nstarted ? d : -1;           // leading zeros carry prev=-1
        total += solve(pos + 1, nprev, tight && (d == cap), nstarted);
    }
    if (!tight) memo[pos][prev + 1][started] = total;
    return total;
}

long long f(long long x) {
    if (x < 0) return 0;
    s = to_string(x);
    memset(memo, -1, sizeof(memo));
    return solve(0, -1, true, false);
}

int main() {
    long long a, b;
    cin >> a >> b;
    cout << f(b) - f(a - 1) << "\n";
    return 0;
}
\end{lstlisting}

\begin{complexitybox}
\textbf{Time Complexity.} $O(\text{digits}\times10\times10)$ per call --
about $20\cdot10\cdot10$ states, each trying $10$ digits: constant for
64-bit inputs. \textbf{Space Complexity.} $O(\text{digits}\times10\times
2)$ for the memo.
\end{complexitybox}

\begin{takeawaybox}
Digit DP counts numbers with a property by building them digit by digit
under the two universal flags \emph{tight} and \emph{started}, memoising
the free (non-tight) states. Range queries become $f(b)-f(a-1)$. This
template adapts to any per-digit constraint by changing the ``skip this
digit'' rule.
\end{takeawaybox}

